\chapter{Reproducibility and artifact evaluation}
\label{ch:repro}

The climax of the note: from the reproducibility crisis to a SWH-native practice for
artifact evaluation. This chapter folds in the standalone \emph{AEC guide} (history of
artifact evaluation, the failures of zip+DOI, the seven properties of Software Heritage, the
organiser/author/reviewer workflow, double-blind via directory \textsc{swhid}s, and the
Spinellis ``hotspots'' strip/restore-qualifiers technique).

\takeaway{The only difference between a double-blind and a camera-ready citation is the
qualifiers; the content hash is identical.}

\section{The reproducibility crisis}
% TODO draft from: reprod-bad-sota.org::#collbergmethod; reproducible-builds.org::#r-b-definition.

\section{Artifact evaluation and badging}
% TODO draft from: sourcecode-reproductibility.org::#artefactevaluation,#artefactevaluation-beyondcs,
% #elife-swh-pipeline,#opensourcereprod. Refs: acmbadging, nisorp31.

\section{Artifact evaluation with Software Heritage}
% TODO: port ../2026-06-AEC-guide/swhid-artifact-evaluation-guide.md to LaTeX here
% (organiser / author / reviewer parts; SWHID verification offline).

\section{Double-blind review and dense pinpointing}
% TODO: dir SWHID for anonymity; Spinellis hotspots. Ref: spinellis2026hotspots.

\begin{handson}[Hands-on: verify a SWHID offline]
\begin{enumerate}[nosep,leftmargin=*]
  \item Given a colleague's file and a claimed content \textsc{swhid}, run
        \texttt{swh identify -{}-verify swh:1:cnt:<hex> file} (exit~0 iff the bytes match).
  \item For a double-blind submission, compute a directory \textsc{swhid}
        (\texttt{swh identify -{}-type directory ./}); it reveals no author or origin, yet
        lets a reviewer confirm the exact bytes.
\end{enumerate}
\end{handson}
